\subsection{The politico-economic equilibrium with LOG preferences}\label{\refsec:PEELOG}
Let $\upsilon_{1,t}^j$ and $\upsilon_{2,t}^j$ denote the indirect utilities of different households (types and generations). Policy is set without commitment, with elections each $t$ aggregating preferences using probabilistic voting that maximizes a weighted sum of indirect utilities. Let $\tau^{t+1}(z_t)$ denote a continuation policy with $z_t$ being a vector of states. The political problem can then be written as
\begin{align*}
    &\max_{\tau_t}\mbox{ }\mathcal{W}_t(z_{t-1}) \equiv \omega\sum_j\left(\gamma_{t-1,j} p_{t-1,j} \mu_{t-1,j} v_{2,t}^j\right)+\nu_t\sum_j\left(\gamma_{t,j}\mu_{t,j} v_{1,t}^j\right) \\
    &\text{s.t. competitive equilibrium and }\tau_{t+1} = \tau^{t+1}(z_t),
\end{align*}
where $\tau^{t+1}(z_t)$ is identified by this process at $t+1$, $\omega$ is a weight reflecting relative political power of the older generation, $\mu_{t,j}$ reflects relative voting shares for different income groups (the old generation at $t$ votes with the voting shares $\mu_{t-1,j}$ they had when young, i.e. shares are attached to a generation, not a period).

With $\rho = 1$, the discount function simplifies to $B_{t+1}^i = \beta_{t,i}$ and the politico-economic equilibrium can be characterized in closed form without any continuation policy; We can exploit that all consumption functions are log-separable in $s_{t-1}$ to show that there are no endogenous states, i.e. the $z_t = \emptyset$ (see the paper for an elaboration of this claim). In general, the first order condition for $\tau_t$ is given by:
\begin{align*}
    \omega\sum_{j}\left(\gamma_{t-1,j}p_{t-1,j}\mu_{t-1,j} \frac{d\upsilon_{2,t}^j}{d\tau_t}\right)+
    \nu_t\sum\left(\gamma_{t,j}\mu_{t,j} \frac{d\upsilon_{1,t}^j}{d\tau_t}\right) = 0,
\end{align*}
where the total derivatives (with "hard d") includes the partial effects of changing policy directly (through economic equilibrium functions), but also the indirect effects working through the continuation policy. 

\begin{subequations}\label{\refeq:PEELOG}
For the log case, however, where there are no effects through the continuation policy, this simplifies significantly. For the consumption smoothers, using the Euler condition, we have $\upsilon_{1,t}^i = (1+\beta_{t,i})\ln\left(\tilde{c}_{1,t}^i\right)+\beta \ln\left(R_{t+1}\right)+\beta \ln(\beta)$. Using that $\tilde{c}_{1,t}^i$ is log-separable in $s_{t-1}$ and the formula for the interest rate, it is straightforward to verify that this reduces to
\begin{align}
    \frac{d\upsilon_{1,t}^i}{d\tau_t} = -\frac{1}{1-\tau_t}\frac{1+\xi}{1+\alpha\xi}\left(1+\beta_{t,i}\frac{\alpha(1+\xi)}{1+\alpha\xi}\right).
\end{align}
For the informal young, the current tax only works through the future informal wage rate:
\begin{align}
    \frac{d\upsilon_{1,t}^0}{d\tau_t} = -\frac{\beta_{t,0}}{1-\tau_t}\frac{\alpha(1+\xi)^2}{(1+\alpha\xi)^2}.
\end{align}
For the older cohorts, the effect works through $\tilde{c}_{2,t}^j$ and, again being log-separable in $s_{t-1}$, this simplifies to
\begin{align}
    \frac{d\upsilon_{2,t}^i}{d\tau_t} =& (1-\alpha)\frac{\partial \ln\left(\Theta_{h,t}\right)}{\partial \tau_t}+\frac{\frac{1-\alpha}{\alpha}\frac{p_{t-1}}{\kappa_{t-1}}\left(1+\theta_t\left[\frac{\eta_{t-1,i}^{1+\xi}}{X_{t-1,i}^{\xi}}\frac{1}{\Gamma_{t-1,h}}-1\right]\right)}{\frac{s_{t-1}^i}{s_{t-1}}+\frac{1-\alpha}{\alpha}\frac{p_{t-1}\tau_t}{\kappa_{t-1}}\left(1+\theta_t\left[\frac{\eta_{t-1,i}^{1+\xi}}{X_{t-1,i}^{\xi}}\frac{1}{\Gamma_{t-1,h}}-1\right]\right)} \\ 
    \frac{d\upsilon_{2,t}^0}{d\tau_t} =& \frac{\dfrac{(1-\alpha)\nu_t\epsilon_t}{\kappa_{t-1}} \Theta_{h,t}^{1-\alpha}\left(1+(1-\alpha)\tau_t\frac{\partial \ln\left(\Theta_{h,t}\right)}{\partial\tau_t}\right)}{\frac{\left(\chi_{t-1}\eta_{t-1,0}\right)^{1+\xi}}{(1+\xi)X_{t-1,0}^{\xi}}\left(\frac{1-\alpha}{\Gamma_{h,t}^{\alpha}}\right)^{\frac{1+\xi}{1+\alpha\xi}}+\dfrac{(1-\alpha)\nu_t\epsilon_t\tau_t}{\kappa_{t-1}} \Theta_{h,t}^{1-\alpha}},
\end{align}
where 
\begin{align*}
    \frac{\partial \ln\left(\Theta_{h,t}\right)}{\partial \tau_t} = -\frac{1}{1-\tau_t}\frac{\xi}{1+\alpha\xi}.
\end{align*}
\end{subequations}





\subsection{The politico-economic equilibrium with CRRA preferences}\label{\refsec:PEE}
With CRRA preferences, the general first order condition is the same, but the total derivatives of indirect utilities change significantly. As we show in \cref{\refnumsec}, we will rely on numerical methods that approximate these derivatives in combination with grid searches. For the young generations, this means that we only have to be able to compute indirect utilities. For the older generations, however, we need to go one step deeper, in order to ensure that we do not need to evaluate on a large grid of the past distribution of savings $s_{t-1,i}/s_{t-1}$ (see the main paper for a thorough explanation). 


\begin{subequations}\label{\refeq:PEE}
We start using the Euler condition for the formal young households, we have that
\begin{align}
    \upsilon_{1,t}^i = (1+B_{t+1}^i)\frac{\left(\tilde{c}_{1,t}^i\right)^{1-1/\rho}}{1-1/\rho}.
\end{align}
The indirect utility of young informal households cannot be reduced further, as they consume "hand-to-mouth", i.e. do not smooth out consumption and utility over time. Thus, we simply have:
\begin{align}
    \upsilon_{1,t}^0 = \frac{1}{1-1/\rho}\left[ \left(\tilde{c}_{1,t}^0\right)^{1-1/\rho}+\beta_{t,0} \left(\tilde{c}_{2,t+1}^0\right)^{1-1/\rho}\right]
\end{align}

Because the policy maker has to take past choices as given (such as savings and relative savings rates), we have to rely on first order conditions to efficiently identify their choices (to avoid searching over a full grid of states $s_{t-1,i}/s_{t-1}$). Whereas we can numerically approximate derivatives of $\upsilon_{1,t}^i$  and $\upsilon_{1,t}^0$ using simple grids, here we have to be a bit more careful. In particular, we have to use:
\begin{align}
    \dfrac{d v_{2,t}^i}{d \tau_t} =& \left(c_{2,t}^i\right)^{1-1/\rho}\dfrac{d\ln\left(c_{2,t}^i\right)}{d\tau_t} \\
    \frac{d\ln\left(c_{2,t}^i\right)}{d\tau_t} =& (1-\alpha)\frac{d \ln\left(h_t\right)}{d \tau_t}+\frac{\frac{1-\alpha}{\alpha}\frac{p_{t-1}}{\kappa_{t-1}}\left(1+\theta_t\left[\frac{\eta_{t-1,i}^{1+\xi}}{X_{t-1,i}^{\xi}}\frac{1}{\Gamma_{t-1,h}}-1\right]\right)}{\frac{s_{t-1}^i}{s_{t-1}}+\frac{1-\alpha}{\alpha}\frac{p_{t-1}\tau_t}{\kappa_{t-1}}\left(1+\theta_t\left[\frac{\eta_{t-1,i}^{1+\xi}}{X_{t-1,i}^{\xi}}\frac{1}{\Gamma_{t-1,h}}-1\right]\right)},
\end{align}
where the second line explicitly uses that $s_{t-1,i}/s_{t-1}$ is predetermined. However, as this ratio has a closed-form expression that relies on $\tau_t$, we do not have to evaluate on the entire grid, rather only on the selection that are actually consistent with economic equilibrium conditions. This is the exact trick that allows us to only use $s_{t-1}$ as a relevant state and not the entire distribution.


Finally, the current old hand-to-mouth consumers have, we do not have to worry about this, we can simply evaluate the indirect utility directly as
\begin{align}
    \upsilon_{2,t}^0 = \frac{\left(\tilde{c}_{2,t}^0\right)^{1-1/\rho}}{1-1/\rho}.
\end{align}
\end{subequations}